\documentclass[11pt]{article}
\usepackage[margin=1in]{geometry}
\usepackage{amsmath,amssymb,amsthm,hyperref}
\newtheorem{theorem}{Theorem}
\newtheorem{lemma}{Lemma}
\newcommand{\fl}[1]{\left\lfloor #1\right\rfloor}
\newcommand{\cl}[1]{\left\lceil #1\right\rceil}
\title{Explicit disjoint zero blocks for square-root floor differences}
\author{Principia Math In Progress}
\date{29 September 2026}
\begin{document}\maketitle
\begin{abstract}
We prove the disjoint zero-block conjecture of Chamberland and Dilcher.
An explicit formula places each block inside the finite interval required
by their alternating-sum problem. The blocks are disjoint and partition
all zero positions. The proof is a bijection obtained by reflecting two
consecutive integers across $n/2$.
\end{abstract}

\section*{Task 1: the original conjecture and its resolution}
Fix an odd positive integer $n$, and put $N=(n-1)/2$. Following
\cite[Equation (2.3) and Conjecture 2.1]{CD}, set
\[
 d_n(\ell)=\fl{\sqrt{2\ell n}}-\fl{\sqrt{(2\ell-1)n}},
 \qquad 1\le\ell\le N.
\]
The source asks that a value $d_n(\lambda)=\delta\ge2$ be associated
with $\delta-1$ consecutive zero positions, with different associations
disjoint. The finite interval is part of the source problem: choosing
zeros arbitrarily far beyond $N$ would not answer it.

\begin{theorem}\label{thm:blocks}
For $1\le\lambda\le N$ with $\delta=d_n(\lambda)\ge2$, define
\[
 a_\lambda=\fl{\sqrt{(2\lambda-1)n}},\qquad
 b_\lambda=\fl{\sqrt{2\lambda n}},\qquad
 L_\lambda=\frac{n+1}{2}+\lambda-b_\lambda.
\]
Then
\[
 B_\lambda=\{L_\lambda,L_\lambda+1,\ldots,L_\lambda+\delta-2\}
 \subseteq\{1,\ldots,N\}
\]
consists entirely of zeros of $d_n$. These blocks are pairwise disjoint,
and their union is exactly $\{\ell:1\le\ell\le N,\ d_n(\ell)=0\}$.
The assertion is vacuous when $n=1$.
\end{theorem}

For example, for $n=33$ the surplus positions are $\lambda=1,2,3$, with
values $3,2,2$. The formula gives respectively $B_1=\{10,11\}$,
$B_2=\{8\}$, and $B_3=\{6\}$, reproducing the source's example.

\section*{Two ceiling identities}
For integer $0\le k\le n$, define $c(k)=\cl{k^2/n}$. Integer translation
inside a ceiling gives the exact reflection identity
\begin{equation}\label{eq:reflection}
 c(n-k)=n-2k+c(k).
\end{equation}
Moreover, for $0\le r\le n-1$,
\begin{equation}\label{eq:gap}
 0\le c(r+1)-c(r)\le2,
\end{equation}
since $((r+1)^2-r^2)/n=(2r+1)/n<2$.
Finally,
\begin{equation}\label{eq:bucket}
 c(k)=2\lambda
 \quad\Longleftrightarrow\quad
 a_\lambda+1\le k\le b_\lambda.
\end{equation}
The strict lower and weak upper square inequalities in this equivalence
also cover the case where an endpoint is a perfect square.

\section*{Constructing each block and proving disjointness}
Fix $\lambda$ as in the theorem. For each
\[
 k\in\{a_\lambda+2,\ldots,b_\lambda\},
 \qquad r=n-k,
 \qquad \ell=N+1+\lambda-k,
\]
equation~\eqref{eq:bucket} gives $c(k-1)=c(k)=2\lambda$.
Applying~\eqref{eq:reflection} to these two integers yields
\begin{equation}\label{eq:two-ceilings}
 c(r)=2\ell-1,\qquad c(r+1)=2\ell+1.
\end{equation}
Therefore
\[
 r^2\le(2\ell-1)n\le2\ell n<(r+1)^2,
\]
and hence
\begin{equation}\label{eq:zero}
 \fl{\sqrt{(2\ell-1)n}}=\fl{\sqrt{2\ell n}}=r.
\end{equation}
Thus $d_n(\ell)=0$.

We also check the interval, rather than assuming it. Since
$\lambda\le N$, we have $b_\lambda\le n-1$, so $r\ge1$.
Consequently $c(r)\ge1$, and~\eqref{eq:two-ceilings} implies $\ell\ge1$.
Also $k-1\ge1$, so $r+1\le n$ and $c(r+1)\le n$;
thus $2\ell+1\le n$ and $\ell\le N$.

As $k$ runs consecutively from $a_\lambda+2$ to $b_\lambda$,
$\ell$ runs consecutively in reverse order through exactly $B_\lambda$.
This proves its size, position, and zero property.
If the same $\ell$ arose from two pairs $(\lambda,k)$, equation
\eqref{eq:zero} would determine the same $r$, and hence the same
$k=n-r$. The equality $\ell=N+1+\lambda-k$ then determines the same
$\lambda$. This proves disjointness.

\section*{The inverse map: every zero occurs}
Let $1\le\ell\le N$ satisfy $d_n(\ell)=0$, and set
\[
 r=\fl{\sqrt{2\ell n}}=\fl{\sqrt{(2\ell-1)n}}.
\]
Then $r^2\le(2\ell-1)n$ and $(r+1)^2>2\ell n$, so
$c(r)\le2\ell-1$ and $c(r+1)\ge2\ell+1$.
Here $1\le r\le n-1$. The gap bound~\eqref{eq:gap} forces equality:
\[
 c(r)=2\ell-1,\qquad c(r+1)=2\ell+1.
\]
Set $k=n-r$ and $\lambda=N+\ell-r$. Reflection now gives
\[
 c(k-1)=c(k)=2\lambda.
\]
In particular $k-1>0$: otherwise the left side would be $c(0)=0$,
while $c(k)>0$. Since $k\le n-1$, the common value is a positive
even integer at most $n-1$; thus $1\le\lambda\le N$.
Equation~\eqref{eq:bucket} shows $a_\lambda+2\le k\le b_\lambda$,
so $d_n(\lambda)\ge2$ and the construction above maps this pair to
the original $\ell$. This proves surjectivity and the theorem.

\paragraph{The already known alternating-sum identity.}
The bijection gives
\[
 \#\{\ell:d_n(\ell)=0\}
   =\sum_{\lambda:d_n(\lambda)\ge2}(d_n(\lambda)-1),
 \qquad \sum_{\ell=1}^{N}d_n(\ell)=N.
\]
It follows that $\sum_{j=1}^{n}(-1)^{j+1}\fl{\sqrt{jn}}=n-N=(n+1)/2$.
This last identity is already Theorem 1.1 of~\cite{CD}; no new credit
is claimed for it. The new target here is the explicit consecutive,
disjoint block structure left conjectural in that paper.

\section*{Tasks 2--5: follow-up questions}
\begin{enumerate}\setcounter{enumi}{1}
\item Determine an asymptotic formula, with an explicit error term,
for the number of zero positions as odd $n\to\infty$.
\item Determine the number of blocks of each length and the maximal
block length, including the dependence on square endpoints.
\item Classify the analogous zero and surplus structure for even $n$,
where reflection preserves, rather than reverses, ceiling parity.
\item Identify higher-power floor differences admitting a comparable
block bijection, and determine which reflection or increment conditions
are necessary for such a construction.
\end{enumerate}

\paragraph{Verification and scope.}
The proof covers every odd positive integer, including composite integers
and square endpoints. The accompanying standard-library verifier uses
integer square roots and integer division only. Its finite exhaustive
and large-integer checks corroborate the proof; they do not replace it.
The source and prior-claim searches are bounded, and their limits are
recorded in the literature audit. No independent peer review is claimed.

\begin{thebibliography}{9}
\bibitem{CD} M.~Chamberland and K.~Dilcher,
\emph{An alternating sum of the floor function of square roots},
\href{https://arxiv.org/abs/2510.26291}{arXiv:2510.26291v1} (2025),
Conjecture 2.1, pp.~2--3. Complete current five-page manuscript checked
29 September 2026.
\bibitem{MathDB} MathDB,
\href{https://mathdb.com/p/371228}{problem 371228},
live status and activity checked 29 September 2026.
\end{thebibliography}
\end{document}
